\begin{recipe}{Pommes Anna Kubusjes}
\kicker[Bijgerecht,Vegetarisch,Frans]{Bijgerecht · vegetarisch · Frans}
\index[register]{Pommes Anna Kubusjes}
\index[register]{Aardappelen}

\meta{40 minuten}

\lettrine{P}{ommes Anna is een klassiek Frans gerecht van dunne aardappelplakjes.} Deze kubusjes-versie geeft dezelfde boterige smaak maar met veel meer krokante randjes. Lekker als bijgerecht bij vlees, vis of ei.

\blockrule{Ingrediënten}
\begin{ingredients}
  \ing{1 kg}{vastkokende aardappelen}
  \ing{100 g}{boter (gesmolten)}
  \ing{2}{teentjes knoflook (geperst)}
  \ing{1 tl}{gedroogde rozemarijn}
  \ingb{zout}
  \ingb{versgemalen zwarte peper}
\end{ingredients}

\blockrule{Bereiding}
\tip{Door ze van tevoren in een vorm te persen worden ze compact en krijgen ze perfecte kubusvorm.}
\begin{steps}
  \item Schil de aardappelen en snijd ze met een mandoline (kaasschaaf of heel scherp mes) in dunne plakjes van ongeveer 2 mm dik. Blancheer ze kort in kokend water tot ze net gaar zijn, giet af en dep ze heel goed droog.
  \item Meng de plakjes in een grote kom met de gesmolten boter, knoflook, zout, peper en rozemarijn. Zorg dat elk plakje goed bedekt is.
  \item Schep de mengeling in een geoliede cakevorm, druk goed aan. Dek af met aluminiumfolie en zet een zwaar gewicht erop (bijvoorbeeld een volle pan of bakplaat met iets zwaar erop).
  \item Zet minstens één nacht in de koelkast om echt compact te worden.
  \item Haal uit de vorm en snijd in nette kubusjes van ongeveer 2-3 cm. Verwarm de oven voor op 225 °C (hete lucht).
  \item Bak 25-35 minuten tot ze goudbruin en krokant zijn. Halverwege omscheppen voor gelijkmatige kleur.
  \item Serveer meteen, eventueel met wat zeezout erover.
\end{steps}
\end{recipe}
